% ============================================================
% Silent Primes and the Variance of the Collision Count
% ============================================================

\documentclass[12pt]{amsart}

\usepackage{amsmath,amssymb,amsthm}
\usepackage{booktabs}
\usepackage{microtype}
\usepackage{url}

\newtheorem{theorem}{Theorem}
\newtheorem{lemma}[theorem]{Lemma}
\newtheorem{proposition}[theorem]{Proposition}
\newtheorem{corollary}[theorem]{Corollary}
\newtheorem{conjecture}[theorem]{Conjecture}
\theoremstyle{definition}
\newtheorem{definition}[theorem]{Definition}
\theoremstyle{remark}
\newtheorem{remark}[theorem]{Remark}

\title{Silent Primes and the Variance\\of the Collision Count}

\author{Alexander S.\ Petty}

\email{alexander.petty@gmail.com}
\date{August 2024}
\subjclass[2020]{11A63, 11N05, 11M41}

\begin{document}
\begin{abstract}
The gate width theorem gives an explicit formula for the
bin-deranging set. This paper develops two consequences.
First, for primitive-root primes at each fixed lag, the
silent primes (those with zero autocorrelation) are the
prime factors of an explicit arithmetic sequence, and
only finitely many exist. Second, computation suggests
the collision count $C(g)$ across non-deranging
multipliers ($p > b{+}1$) has variance proportional to
$Q = \lfloor (p{-}1)/b \rfloor$ (conjectured, not
proved).
\end{abstract}
\maketitle

% ============================================================
\section{Introduction}

This paper studies the variance of the collision count and introduces the Dirichlet series $L^G(s)$.

The gate width theorem~\cite{paper11} proved that the digit
function $\delta(r) = \lfloor br/p \rfloor$ has exactly $b - 1$
bin-deranging elements for any prime $p > b$, and identified them
as the M\"obius family $g = -u/(b{-}u) \pmod{p}$ for
$u = 1, \ldots, b{-}1$. The present paper asks two questions.
First: as $p$ varies, which primes are silent at a given lag?
Second: how does the collision count fluctuate across the
non-deranging multipliers?

The first question leads to an explicit arithmetic
characterization of the silent primes. The second reveals a
square-root variance law that connects the fractional field to
the equidistribution theory of sequences modulo~$p$.  The
Dirichlet series framework follows
classical methods~\cite{davenport,hardy}.

% ============================================================
\section{Silent Primes}

\begin{definition}
A prime $p > b$ is \emph{silent at lag~$\ell$} if $R_p(\ell) = 0$.
\end{definition}

For a prime $p$ where $b$ is a primitive root, $R_p(\ell) =
C(b^{\ell} \bmod p)$. By the gate width theorem, $C(g) = 0$ if
and only if $g \equiv -u/(b{-}u) \pmod{p}$ for some
$u \in \{1, \ldots, b{-}1\}$.

\begin{theorem}\label{thm:silent}
For a primitive-root prime $p > b$ at fixed lag~$\ell \ge 1$,
$p$ is silent at lag~$\ell$ (i.e., $R_p(\ell) = 0$) if and
only if $p$ divides
\[
N(\ell, u) = b^{\ell+1} + u(1 - b^{\ell})
\]
for some $u \in \{1, \ldots, b{-}1\}$. Only finitely many
primes are silent at any given lag.
\end{theorem}

\begin{proof}
The condition $b^{\ell} \equiv -u/(b{-}u) \pmod{p}$ gives
$b^{\ell}(b{-}u) + u \equiv 0 \pmod{p}$, i.e.,
$p \mid N(\ell, u)$. Since $N(\ell, u) \ne 0$ for $\ell \ge 1$,
only finitely many primes divide it.
\end{proof}

\begin{proposition}\label{prop:arithmetic}
The $b - 1$ integers $N(\ell, u)$ for $u = 1, \ldots, b{-}1$
form an arithmetic sequence with first term
$b^{\ell}(b{-}1) + 1$ and common difference $-(b^{\ell} - 1)$.
\end{proposition}

\begin{proof}
$N(\ell, u) = b^{\ell+1} - u(b^{\ell} - 1)$. The common
difference is $-(b^{\ell} - 1)$.
\end{proof}

\begin{table}[h]
\centering
\begin{tabular}{rrl}
\toprule
$\ell$ & silent primes ($p > 10$) & from $N(\ell, u)$ \\
\midrule
$1$ & $11, 13, 19, 23, 37, 41, 73$ & factors of
  $91, 82, \ldots, 19$ \\
$2$ & $13, 17, 19, 29, 37, 53, 101, \ldots$ & factors of
  $901, 802, \ldots, 109$ \\
$3$ & $11, 13, 19, 31, 47, 79, 97, \ldots$ & factors of
  $9001, \ldots, 1009$ \\
\bottomrule
\end{tabular}
\caption{Silent primes at the first three lags in base~$10$.}
\label{tab:silent}
\end{table}

% ============================================================
\section{The Collision Parameter and the Sawtooth Form}

For a non-deranging $g \ne 1$, the gate width proof introduces
the collision parameter $c = b(1{-}g)^{-1} \pmod{p}$ with
$c \in \{b{+}1, \ldots, p{-}1\}$.

Write $\left(\!\left( x \right)\!\right) = \{x\} - \tfrac{1}{2}$
for the sawtooth function.

\begin{theorem}[Sawtooth form]\label{thm:sawtooth}
With the sawtooth notation above,
\[
C(g) = Q + \sum_{m=1}^{Q}
\operatorname{sgn}\!\left(
\left(\!\left(\frac{mc}{p}\right)\!\right) -
\left(\!\left(\frac{mb}{p}\right)\!\right)
\right),
\]
where $Q = \lfloor (p{-}1)/b \rfloor$ and
$c = b(1 - g)^{-1} \bmod p$ is the collision parameter.
\end{theorem}

\begin{proof}
The collision condition at residue $r$ with parameter $c$
reduces to a comparison of sawtooth values at
$m = 1, \ldots, Q$. For each $m$: if
$((mc/p)) > ((mb/p))$, both $r$ and $p - r$ produce
collisions (by the bilateral symmetry $r \mapsto p - r$,
which preserves the comparison). So each positive
comparison contributes $2$ collisions. If
$((mc/p)) < ((mb/p))$, neither $r$ nor $p - r$
collides, contributing $0$. Writing
$\mathbf{1}_{>} = (1 + \operatorname{sgn})/2$ and
doubling: the total collision count is
$2 \cdot \sum_{m=1}^{Q} \frac{1 + \operatorname{sgn}}{2}
= Q + \sum \operatorname{sgn}$.
\end{proof}

The sawtooth form expresses the collision count as a signed
comparison of two sequences modulo~$p$: the sequence $mc \bmod p$
(controlled by the collision parameter) and the reference
sequence $mb$ (controlled by the base). The collision count
measures how often the first sequence ``leads'' the second.

The collision count depends on $c$ through the wrapping
behavior of the sequence $mc \bmod p$. The exact
relationship between $c$ and $C(g)$ is governed by the
sawtooth comparison above.

% ============================================================
\section{The Variance Conjecture}

The collision count $C(g)$ varies across the non-deranging
multipliers. We now characterize its fluctuations.

\begin{definition}
For a prime $p > b$, the \emph{collision variance} is
\[
V(p) = \frac{1}{|S|} \sum_{g \in S} \bigl(C(g) -
\overline{C}\bigr)^2,
\]
where $S = \{g \in \mathbb{F}_p^{\times} : C(g) > 0,\,
g \ne 1\}$ and $\overline{C}$ is the mean of $C$ over~$S$.
\end{definition}

\begin{conjecture}[Variance scaling]\label{conj:variance}
The collision variance satisfies
\[
V(p) = \lambda_b \cdot Q + O(1),
\]
where $\lambda_b > 0$ is a constant depending only on the
base~$b$, and $Q = \lfloor (p{-}1)/b \rfloor$.

Equivalently, the standard deviation of $C(g)$ over the
non-deranging multipliers is $O(\sqrt{Q}) = O(\sqrt{p/b})$.
\end{conjecture}

This conjecture has been verified computationally for all primes
$p \le 10000$ in bases $b \in \{3, 6, 7, 10, 12\}$. The
constant $\lambda_b$ depends on $b$ and converges slowly as
$Q \to \infty$. In base~$10$: $\lambda_{10} \approx 0.89$.

\begin{table}[h]
\centering
\begin{tabular}{rrrrr}
\toprule
$p$ & $Q$ & $V(p)$ & $V/Q$ & $\sigma/\!\sqrt{Q}$ \\
\midrule
$97$ & $9$ & $6.00$ & $0.667$ & $0.817$ \\
$193$ & $19$ & $13.97$ & $0.735$ & $0.857$ \\
$499$ & $49$ & $40.45$ & $0.826$ & $0.909$ \\
$997$ & $99$ & $84.35$ & $0.852$ & $0.923$ \\
$1999$ & $199$ & $173.10$ & $0.870$ & $0.933$ \\
$4999$ & $499$ & $440.55$ & $0.883$ & $0.940$ \\
$9973$ & $997$ & $884.83$ & $0.888$ & $0.942$ \\
\bottomrule
\end{tabular}
\caption{Collision variance in base~$10$. The ratio $V/Q$
converges to $\lambda_{10} \approx 0.89$, giving
$\sigma/\!\sqrt{Q} \to 0.94$.}
\label{tab:variance}
\end{table}

\begin{remark}
The square-root scaling $\sigma(C) \sim \sqrt{Q}$ is the
characteristic fluctuation of equidistributed sequences
modulo~$p$. The sawtooth comparison in
Theorem~\ref{thm:sawtooth} sums $Q$ terms, each $\pm 1$, whose
signs depend on the equidistribution of $mc \bmod p$ relative to
$mb$. When the signs are approximately independent (as occurs for
generic $c$), the central limit theorem gives variance
proportional to $Q$. This scaling is observed computationally;
a rigorous proof that the signs are sufficiently independent
remains open.

The collision count $C(g)$ therefore fluctuates around its mean
$\overline{C} \approx Q$ with amplitude $\sqrt{Q}$. This is
the same square-root fluctuation pattern that governs the error
term in the prime number theorem under the Riemann Hypothesis:
$\pi(x) - \operatorname{li}(x) = O(\sqrt{x} \log x)$.
Whether the fluctuations of $C(g)$ and the fluctuations of
$\pi(x)$ are controlled by the same mechanism (the zeros of an
$L$-function) is the central open question.
\end{remark}

% ============================================================
\section{The Dirichlet Series}

\begin{definition}
The \emph{fractional field Dirichlet series} at lag~$\ell$ is
\[
L^{G}(s, \ell) = \sum_{p > b} \frac{R_p(\ell)}{p^s}.
\]
\end{definition}

\begin{theorem}\label{thm:convergence}
$L^{G}(s, \ell)$ converges absolutely for
$\operatorname{Re}(s) > 2$.
\end{theorem}

\begin{proof}
The trivial bound $0 \le R_p(\ell) \le p - 1$ gives
$|R_p(\ell) / p^s| \le p^{1-\operatorname{Re}(s)}$.
The sum $\sum_p p^{1-s}$ converges for
$\operatorname{Re}(s) > 2$.
\end{proof}

At each lag~$\ell$, the series runs over primes $p > b$.
For primitive-root primes,
$R_p(\ell) = C(b^{\ell} \bmod p)$, connecting the
autocorrelation to the global collision count. For
non-primitive-root primes, this equality does not hold
in general.

If the variance conjecture holds, the individual terms
fluctuate by $O(\sqrt{p/b})$ around the mean. Whether
these fluctuations have additional analytic structure
(meromorphic continuation, connection to $L$-functions)
is an open question.

% ============================================================
\section{The M\"obius Family}

The deranging set $\{-u/(b{-}u) : u = 1, \ldots, b{-}1\}$ is a
family of $b - 1$ rational numbers, independent of~$p$.
In base~$10$:
\[
-\tfrac{1}{9},\quad -\tfrac{1}{4},\quad -\tfrac{3}{7},\quad
-\tfrac{2}{3},\quad -1,\quad -\tfrac{3}{2},\quad
-\tfrac{7}{3},\quad -4,\quad -9.
\]

\begin{proposition}\label{prop:mobius-structure}
The M\"obius family satisfies:
\begin{enumerate}
\item \textbf{Inversion symmetry.} $-u/(b{-}u)$ and
      $-(b{-}u)/u$ are multiplicative inverses. The deranging
      set is closed under inversion.
\item \textbf{Complement element.} For even $b$, $u = b/2$
      gives $g \equiv -1 \pmod{p}$.
\item \textbf{Complementary structure.} The fractions $u/(b{-}u)$ for
      $u = 1, \ldots, b{-}1$ are the ratios of complementary
      parts of~$b$.
\end{enumerate}
\end{proposition}

The reduction of the M\"obius family modulo varying primes
governs both the silent prime structure (Section~2) and the
collision variance (Section~4). At each lag~$\ell$, the silent
primes are those where $b^{\ell}$ reduces to a member of the
family. The collision variance measures how the non-family
elements distribute their collision counts.

% ============================================================
\section{Remarks}

\subsection*{Toward the fluctuation sum}

The decomposition of $L^{G}(s, \ell)$ into mean and fluctuation
parts sets up the central question for the next paper: does the
fluctuation sum
\[
\Phi_x(s, \ell) = \sum_{b+1 < p \le x}
\frac{C(b^{\ell} \bmod p) - \overline{C}_p}{p^s}
\]
have meromorphic continuation beyond its half-plane of
convergence? Assuming the variance conjecture
(Conjecture~\ref{conj:variance}), the individual
fluctuation terms would be $O(\sqrt{p/b})$, suggesting
(but not proving) convergence for
$\operatorname{Re}(s) > 3/2$.

If $F$ admits continuation to $\operatorname{Re}(s) > 1$ with
controlled growth, the zeros of $F$ in the critical strip would
govern the distribution of collision fluctuations across primes,
in direct analogy to how the zeros of $\zeta(s)$ govern the
distribution of primes.

\subsection*{The variance constant}

The constant $\lambda_b$ in Conjecture~\ref{conj:variance} depends
on~$b$ and converges slowly with~$Q$. Its exact value is an open
problem. The sawtooth form (Theorem~\ref{thm:sawtooth}) suggests
that $\lambda_b$ is determined by the equidistribution properties
of the sequence $mc \bmod p$ as $c$ varies over
$\{b{+}1, \ldots, p{-}1\}$, connecting the variance to the
discrepancy theory of modular sequences.  The logarithmic
divergence of $\sum 1/p$ that underpins the Mertens-type
growth law was established by Mertens~\cite{mertens}; see
also~\cite{hardy}.

% ============================================================
\begin{thebibliography}{9}

\bibitem{paper11}
A.~S.~Petty,
Bin derangements and the gate width theorem,
research note, 2024.

\bibitem{davenport}
H.~Davenport,
\emph{Multiplicative Number Theory},
3rd~ed., Springer, 2000.

\bibitem{hardy}
G.~H.~Hardy and E.~M.~Wright,
\emph{An Introduction to the Theory of Numbers},
6th~ed., Oxford University Press, 2008.

\bibitem{mertens}
F.~Mertens,
Ein Beitrag zur analytischen Zahlentheorie,
\emph{J.\ Reine Angew.\ Math.}\ \textbf{78} (1874), 46--62.

\end{thebibliography}

\end{document}
